\chapter{Arrêt optimal}\label{ch:iv1:optimal-stopping}

\begin{ivtrains}
Décider quand s'arrêter, en raisonnant à rebours à partir de la dernière décision. Un seuil est la valeur de la poursuite, et la plupart de ces questions se ramènent à la détermination de ce nombre. Les dernières demandent ce que vaut une information, et quand aucune règle d'arrêt ne peut aider.
\end{ivtrains}

\begin{question}[$\star$ \iqroles{QR, QT, Trader, Analyste} \iqfirm{toutes} \iqlevel{junior}]\label{iq:iv1:optimal-stopping:1}
Vous lancez un dé équilibré et vous recevez en dollars sa valeur, mais vous pouvez refuser le lancer et relancer le dé en prenant la nouvelle valeur. Avec une seule relance possible, quand gardez-vous le résultat, et combien vaut le jeu~? Et avec deux relances possibles~?
\end{question}

\begin{question}[$\star$ \iqroles{QR, QT, Trader} \iqfirm{toutes} \iqlevel{junior}]\label{iq:iv1:optimal-stopping:2}
Vous tirez un nombre uniforme sur $[0, 1]$ et vous pouvez le garder ou le rejeter pour un nouveau tirage. Combien vaut le jeu avec une seule possibilité de rejet~? Avec deux~? Avec trois~? Quel est le seuil à chaque étape~?
\end{question}

\begin{question}[$\star$ \iqroles{QR, QT, Trader} \iqfirm{toutes} \iqlevel{junior}]\label{iq:iv1:optimal-stopping:3}
Un sac contient 100 tickets, numérotés de 1 à 100. Vous en tirez un et vous le voyez. Vous pouvez le garder et recevoir en dollars le numéro qui y figure, ou tirer un second ticket parmi les 99 restants et recevoir à la place le numéro de celui-ci. Que devez-vous faire, et combien vaut le jeu~?
\end{question}

\begin{question}[$\star$ \iqroles{QR, QT, Trader, Analyste} \iqfirm{toutes} \iqlevel{junior}]\label{iq:iv1:optimal-stopping:4}
Vous retournez les cartes d'un jeu standard mélangé, une à une. À tout moment, vous pouvez dire «~stop~», et vous gagnez si la carte suivante est rouge~; si vous ne dites jamais rien, vous êtes obligé de parier sur la dernière carte. Quelle règle vous donne la meilleure chance de gagner~?
\end{question}

\begin{question}[$\star\star$ \iqroles{QR, QT, Trader, DS} \iqfirm{toutes} \iqlevel{junior}]\label{iq:iv1:optimal-stopping:5}
Vous interrogez $n$ candidats dans un ordre aléatoire. Après chaque entretien, vous ne pouvez comparer le candidat qu'avec ceux déjà rencontrés, et vous devez soit l'embaucher sur-le-champ, soit le rejeter définitivement. Vous ne gagnez que si vous embauchez le meilleur des $n$. Trouvez la meilleure règle et la probabilité de gagner pour $n = 3$, puis la meilleure règle pour $n = 10$, puis la limite pour $n$ grand.
\end{question}

\begin{question}[$\star\star$ \iqroles{QR, QT, Trader, Analyste} \iqfirm{toutes} \iqlevel{junior}]\label{iq:iv1:optimal-stopping:6}
Vous pouvez tirer des nombres uniformes sur $[0, 1]$ autant de fois que vous le souhaitez, en payant $0.02$ pour chaque tirage, le premier compris, et vous gardez le dernier nombre tiré. Quand vous arrêtez, combien de tirages effectuez-vous en moyenne, et combien vaut le jeu~? Que se passe-t-il quand le coût augmente, et quand il tend vers 0~?
\end{question}

\begin{question}[$\star\star$ \iqroles{QR, QT, Trader} \iqfirm{toutes} \iqlevel{junior}]\label{iq:iv1:optimal-stopping:7}
Vous lancez gratuitement un dé équilibré à 20 faces, et vous recevez le nombre affiché. Vous pouvez relancer autant de fois que vous le voulez, en payant 1~\$ à chaque relance, et seul le dernier nombre compte. Quand vous arrêtez, et combien vaut le jeu~?
\end{question}

\begin{question}[$\star\star$ \iqroles{QR, QT, Analyste, DS} \iqfirm{toutes} \iqlevel{junior}]\label{iq:iv1:optimal-stopping:8}
Vous vendez un bloc d'actions. Chaque jour, une offre arrive, uniforme sur $[0, 1]$, et une offre refusée est perdue. Une somme reçue un jour plus tard vaut $0.8$ fois sa valeur actuelle. Quelles offres acceptez-vous, combien de temps faut-il en moyenne, et combien vaut le bloc juste avant l'arrivée de la première offre~?
\end{question}

\begin{question}[$\star\star$ \iqroles{QR, QT, Trader} \iqfirm{toutes} \iqlevel{junior}]\label{iq:iv1:optimal-stopping:9}
Vous lancez un dé équilibré à plusieurs reprises et additionnez les résultats. Vous pouvez vous arrêter à tout moment et prendre le total en dollars, mais si vous obtenez un 6, le jeu s'arrête et vous ne recevez rien. Quand devez-vous vous arrêter~? Environ combien vaut le jeu~?
\end{question}

\begin{question}[$\star\star$ \iqroles{QR, QT, Trader} \iqfirm{toutes} \iqlevel{junior}]\label{iq:iv1:optimal-stopping:10}
Une pièce équilibrée est lancée jusqu'à 4 fois, et vous pouvez vous arrêter après n'importe quel lancer, y compris avant le premier. Vous recevez le nombre de piles moins le nombre de faces obtenus jusqu'alors, ce qui peut être négatif. Quelle est la meilleure règle, et combien vaut-elle~? Comment la réponse change-t-elle si vous recevez le nombre de piles moins le nombre de faces ou zéro, selon le plus grand des deux~?
\end{question}

\begin{question}[$\star\star$ \iqroles{QR, QT, Trader, Analyste} \iqfirm{hedge fund} \iqlevel{junior}]\label{iq:iv1:optimal-stopping:11}
Une transaction coûte 40 et rapporte 100 si un deal se conclut, ce que vous estimez avoir une probabilité de 50~\% de se produire, et rien sinon. Un rapport de recherche dira «~oui~» ou «~non~», et il a raison 80~\% du temps, quelle que soit la réalité. Quel est le montant maximal que vous paieriez pour ce rapport~? Que paieriez-vous pour un rapport parfait~?
\end{question}

\begin{question}[$\star\star$ \iqroles{QR, QT, DS, Trader} \iqfirm{toutes} \iqlevel{junior}]\label{iq:iv1:optimal-stopping:12}
La pièce A a une probabilité connue de face de $1/2$, et la pièce B a une probabilité inconnue de face, uniforme sur $[0, 1]$. Vous lancez une pièce à la fois, en choisissant chaque fois après avoir vu les résultats passés, et vous recevez 1~\$ par face. Avec deux lancers, laquelle lancez-vous en premier, et combien vaut le jeu~? Pourquoi lancer une pièce qui a la même probabilité moyenne que l'autre~?
\end{question}

\begin{question}[$\star\star$ \iqroles{QR, QT, Trader} \iqfirm{teneur de marché} \iqlevel{junior}]\label{iq:iv1:optimal-stopping:13}
Vous tirez des nombres uniformes sur $[0, 1]$ et les additionnez. Vous pouvez vous arrêter à tout moment et prendre le total, mais si ce total dépasse 1, vous ne recevez rien. Quand devez-vous vous arrêter, et combien vaut le jeu au départ~?
\end{question}

\begin{question}[$\star\star\star$ \iqroles{QR, QT, Trader} \iqfirm{teneur de marché, hedge fund} \iqlevel{junior, confirmé}]\label{iq:iv1:optimal-stopping:14}
Un jeu de cartes est distribué face visible, une carte après l'autre. Chaque carte rouge vous rapporte 1~\$, chaque carte noire vous coûte 1~\$, et vous pouvez vous arrêter à tout moment. Calculez à la main la valeur d'un jeu de 2 cartes rouges et 2 noires, et décrivez la méthode pour un jeu complet de 26 et 26. Environ combien vaut le jeu complet~?
\end{question}

\begin{question}[$\star\star\star$ \iqroles{QR, QT, Trader} \iqfirm{hedge fund} \iqlevel{junior, confirmé}]\label{iq:iv1:optimal-stopping:15}
Vous conduisez le long d'une longue rue droite en direction d'un restaurant. Des places de stationnement sont espacées d'une unité sur toute la rue, l'une d'elles se trouvant exactement devant le restaurant, et chacune est libre avec une probabilité de $0.1$, indépendamment des autres. Vous ne voyez si une place est libre qu'en l'atteignant, et vous ne pouvez pas faire demi-tour. Si vous vous garez à $x$ unités du restaurant, avant ou après lui, vous marchez $x$ unités. Quelles places libres prenez-vous, et quelle est la distance moyenne à parcourir à pied~? Quelle est la règle pour une probabilité générale $p$ qu'une place soit libre~?
\end{question}

\begin{question}[$\star\star\star$ \iqroles{QR, QT, DS, Trader} \iqfirm{hedge fund} \iqlevel{junior, confirmé}]\label{iq:iv1:optimal-stopping:16}
Une pièce a une probabilité inconnue de face, uniforme sur $[0, 1]$. Vous placerez un pari sur le prochain lancer, à cote égale et pour 1~\$, du côté de votre choix. Auparavant, vous pouvez observer $n$ lancers gratuits. Quel est votre gain espéré pour $n = 0$, 1, 2 et 3~? Qu'apporte le second lancer, et quelle est la limite lorsque $n$ croît~?
\end{question}

\begin{question}[$\star\star\star$ \iqroles{QR, QT, Trader} \iqfirm{teneur de marché, hedge fund} \iqlevel{junior, confirmé}]\label{iq:iv1:optimal-stopping:17}
Le joueur A tire un nombre uniforme sur $[0, 1]$ et peut le rejeter une fois pour un nouveau tirage, qu'il doit alors garder. Le joueur B voit le nombre final de A, puis fait de même, avec un rejet possible. Le plus grand nombre gagne. Quand A doit-il garder le premier tirage, et quelle est sa probabilité de gagner~?
\end{question}

\begin{question}[$\star\star\star$ \iqroles{QR, QT, Trader, DS} \iqfirm{hedge fund} \iqlevel{junior, confirmé}]\label{iq:iv1:optimal-stopping:18}
Il y a 30 essais indépendants dans l'ordre, chacun étant un succès avec une probabilité de $0.1$. Vous les observez un à un, et vous pouvez vous arrêter une seule fois, sur un succès. Vous gagnez si ce succès s'avère être le dernier des 30. Quelle règle est la meilleure, et quelle est la probabilité de gagner~? Comment ce problème se relie-t-il à l'embauche du meilleur de $n$ candidats~?
\end{question}

\begin{question}[$\star\star\star\star$ \iqroles{QR, QT} \iqfirm{hedge fund, teneur de marché} \iqlevel{confirmé}]\label{iq:iv1:optimal-stopping:19}
Vous lancez une pièce équilibrée autant de fois que vous le voulez, et quand vous arrêtez, vous recevez la fraction des lancers qui ont donné face. Il n'y a ni coût ni limite. Montrez que le jeu vaut au moins $3/4$. Déterminez ce que vaut, sous forme close, la règle «~arrêter la première fois que les piles mènent~». Vous arrêteriez-vous jamais avec autant de piles que de faces, ou même en retard~?
\end{question}

\begin{question}[$\star\star\star\star$ \iqroles{QR, QT} \iqfirm{hedge fund} \iqlevel{confirmé}]\label{iq:iv1:optimal-stopping:20}
Des gains non négatifs indépendants $X_1, \dots, X_n$, de lois connues, sont présentés dans l'ordre. Vous pouvez prendre un gain au moment où il est présenté, ou le perdre définitivement. Un prophète, qui les voit tous à l'avance, obtient $E[\max X_i]$. Montrez qu'un seuil fixe unique vous permet d'obtenir au moins la moitié de cette somme, quelles que soient les lois, et montrez par un exemple que vous ne pouvez pas garantir plus de la moitié.
\end{question}
